% !TEX program = xelatex
\documentclass[a4paper,10pt]{article}

% ─── Packages ───
\usepackage[margin=0.5in, top=0.35in, bottom=0.3in]{geometry}
\usepackage{fontspec}
\usepackage{xeCJK}
\usepackage{titlesec}
\usepackage{enumitem}
\usepackage{hyperref}
\usepackage{xcolor}

% ─── Fonts: auto-detect ───
\setmainfont{TeX Gyre Termes}

\IfFontExistsTF{Noto Serif CJK SC}{
  \setCJKmainfont{Noto Serif CJK SC}[BoldFont=Noto Sans CJK SC Bold]
  \setCJKsansfont{Noto Sans CJK SC}
}{
  \IfFontExistsTF{SimSun}{
    \setCJKmainfont{SimSun}[BoldFont=SimHei]
    \setCJKsansfont{SimHei}
  }{
    \setCJKmainfont{FandolSong}[BoldFont=FandolHei]
    \setCJKsansfont{FandolHei}
  }
}

% ─── Colors ───
\definecolor{accent}{HTML}{2B5A7C}
\definecolor{lightgray}{HTML}{888888}

% ─── Hyperlinks ───
\hypersetup{colorlinks=true, urlcolor=accent, linkcolor=accent}

% ─── Section formatting ───
\titleformat{\section}{\large\bfseries\color{accent}}{}{0em}{}[\vspace{-6pt}\textcolor{accent}{\rule{\textwidth}{0.4pt}}]
\titlespacing{\section}{0pt}{8pt}{4pt}

% ─── Custom commands ───
\newcommand{\cvheader}[4]{%
  \begin{center}
    {\LARGE\bfseries #1 \; \normalsize\color{lightgray}#2}\\[4pt]
    {\small #3 \quad $\cdot$ \quad \href{#4}{github.com/YueFangfly} \quad $\cdot$ \quad 金华}
  \end{center}
  \vspace{-4pt}
}

\newcommand{\projectentry}[4]{%
  \vspace{1pt}
  \noindent\textbf{#1} \hfill {\small\color{lightgray}#2}\\[-3pt]
  {\small\color{accent}#3} \hfill {\small #4}
  \vspace{0pt}
}

\newcommand{\eduentry}[4]{%
  \noindent\textbf{#1} \hfill {\small\color{lightgray}#2}\\[-2pt]
  {\small #3} \hfill {\small #4}
}

% ─── List style ───
\setlist[itemize]{left=0.8em, itemsep=1pt, parsep=0pt, topsep=2pt}

% ─── Reduce spacing ───
\setlength{\parindent}{0pt}
\setlength{\parskip}{0pt}
\pagestyle{empty}

\begin{document}

% ══════════════════════════════════════════
% HEADER
% ══════════════════════════════════════════

\cvheader{Yue Fang}{方玥}{fangyuefly@qq.com}{https://github.com/YueFangfly}

\vspace{-6pt}
\begin{center}
{\small 统计学背景，具备 AI 产品从 0 到 1 的独立交付能力。擅长将 LLM 能力转化为用户可感知的产品功能，熟悉 RAG、Agent 架构与全栈开发。}
\end{center}
\vspace{-4pt}

% ══════════════════════════════════════════
% EDUCATION
% ══════════════════════════════════════════

\section*{教育背景}

\eduentry{University of Nottingham 诺丁汉大学}{2022 -- 2026}{BSc Statistics 统计学学士（2+2：宁波诺丁汉 + 英国校区）}{}

\vspace{2pt}
{\small\color{lightgray}\textbf{核心课程：}\color{black}概率论与数理统计、随机过程、时间序列分析、回归分析、多元统计、数值方法（PDE/SDE）、统计计算（R \& Python）}

% ══════════════════════════════════════════
% SKILLS
% ══════════════════════════════════════════

\section*{技术能力}

\vspace{1pt}
{\small
\textbf{编程语言：} Python, Swift, R, MATLAB, SQL, C++ (Arduino) \quad
\textbf{AI/LLM：} Claude API (Tool Use), OpenAI API, DeepSeek API, 提示词工程 \\[2pt]
\textbf{RAG/Agent：} RAG (检索增强生成), ReAct Agent, ChromaDB, Embedding (all-MiniLM-L6-v2), 语义分块, 向量检索 \\[2pt]
\textbf{后端/部署：} FastAPI, Pydantic, Docker, RESTful API, Uvicorn \quad
\textbf{iOS开发：} SwiftUI, HealthKit, EventKit, Speech框架, SQLite \\[2pt]
\textbf{数据分析：} SARIMA, Monte Carlo, PDE/SDE数值方法, 离散事件仿真 \quad
\textbf{工具：} Git, GitHub, Xcode, Arduino IDE, Firebase
}

% ══════════════════════════════════════════
% PROJECTS
% ══════════════════════════════════════════

\section*{项目经历}

% ─── Keke App ───
\projectentry{Keke AI 陪伴 App（iOS + Python 后端）}{独立项目}{产品设计 \& 全栈开发}{\href{https://github.com/YueFangfly/Keke-AI-accompany-}{GitHub $\nearrow$}}

\begin{itemize}
  \item 独立完成产品构思与全流程开发，iOS 原生 App（57 个 Swift 文件）+ Python 后端，支持 Claude / GPT / DeepSeek / Grok 四家大模型
  \item \textbf{RAG}：语义分块 $\rightarrow$ Embedding（all-MiniLM-L6-v2）$\rightarrow$ ChromaDB 向量存储 $\rightarrow$ Cosine Similarity 检索 $\rightarrow$ 上下文注入 LLM，实现长期记忆精准召回
  \item \textbf{ReAct Agent}：多轮 Tool Calling 循环（最多 10 轮），编排 7 个工具（联网搜索、记忆读写、天气、健康分析等），对话后自动提取关键信息存入向量库
  \item \textbf{后端}：FastAPI + Pydantic + Docker 容器化，RESTful API（/chat、/memory/search 等），22 个单元测试全覆盖
  \item 记忆系统：SQLite + FTS5 全文检索 + trigram 分词 + BM25 排序 + 情绪坐标（valence/arousal）上下文感知检索
  \item 语音通话（Speech Recognition + ElevenLabs TTS）、LLM Tool Use 操控系统提醒/日历、6 维情绪模型、HealthKit 集成
\end{itemize}

% ─── DPDP ───
\projectentry{动态配送调度优化（DPDP）}{小组项目}{数学建模与仿真}{2025}

\begin{itemize}
  \item 建模\textbf{动态取送货问题}（DPDP），扩展 TSP/VRP 至实时动态订单环境，非齐次 Poisson 过程模拟时变到达（$\lambda(t)$: 0.3--1.5 单/分钟），离散事件仿真引擎驱动
  \item 对比 Nearest Driver（贪心）与 Just-in-Time（ACA 预测同步）两种调度策略；构建随机行驶时间模型（路网修正 + 拥堵系数 + 有界正态延迟），灵敏度分析确定最优车队规模（21--25 名）
\end{itemize}

% ─── Evacuation ───
\projectentry{人群疏散建模：PDE 与 Agent-Based 方法比较}{毕业论文（小组项目）}{建模与仿真开发}{2026}

\begin{itemize}
  \item 对比 PDE 模型（Hughes 连续性方程 + 有限体积法）与 Agent-Based 模型（SDE + Euler-Maruyama + Brownian Motion），100 次 Monte Carlo 模拟
  \item 出口布局优化减少 23--35\% 疏散时间；M/Er/2 排队模型模拟瓶颈动力学；人群异质性建模（$v_{\max} = 0.75$ m/s），对标英国 150 秒法规
\end{itemize}

% ─── Time Series ───
\projectentry{英格兰温度时间序列分析}{个人课程作业}{数据分析}{2025}

\begin{itemize}
  \item 使用 R 对 75 年月度温度数据建模，ADF 检验 + ACF/PACF 诊断确定 SARIMA(0,1,0)(1,0,1)$_{12}$ 阶数，Ljung-Box 验证残差白噪声，实现样本外预测
\end{itemize}

% ─── IoT ───
\projectentry{IoT 陪伴设备}{个人项目 · 进行中}{硬件 \& 固件开发}{}

\begin{itemize}
  \item 基于 ESP32 设计物联网陪伴设备，集成触摸/温度/运动传感器；WiFi 云中继（Firebase）+ BLE 本地协议双连接方案
\end{itemize}

% ══════════════════════════════════════════
% LANGUAGES
% ══════════════════════════════════════════

\section*{语言能力}
\vspace{-2pt}
{\small \textbf{中文（普通话）} 母语 \qquad \textbf{英语} 流利（2+2英国学位项目）}

\end{document}
